% !TEX root = ../main.tex
% !TEX program = lualatex

\documentclass[../main.tex]{subfiles}
\externaldocument{../main}

\begin{document}

\section{Manufacturing and Producibility}

\subsection{Why manufacturing belongs in acquisition decisions}

Manufacturing is not a post-design activity.  A program can meet technical
requirements on paper and still fail because the design cannot be produced,
tested, funded, or sustained at the required rate.  Coursebook 3.7.1 frames the
manufacturing decision as a cost, schedule, and performance control problem:
the program manager must use producibility thinking early enough to influence
system design, supplier choices, manufacturing planning, and quality management
before production commitments become expensive to reverse
\autocite{edo-3-7-1-manufacturing-producibility-2025}.

\subsection{Nine industrial process elements}

\begin{longtblr}[
  caption = {Industrial process elements to check before production commitment.},
  label = {tab:industrial-process-elements},
]{
  width = \linewidth,
  colspec = {X[1.05,l] X[0.9,l] X[1.65,l]},
  rowhead = 1,
  hlines,
  vlines,
}
\textbf{Element} & \textbf{Program-manager leverage} & \textbf{Board cue} \\
Design & Direct influence & Push producibility, testability, modularity, tolerances, and design margins into the design baseline. \\
Materials & Direct influence & Check critical materials, alternate sources, lead times, and obsolescence risk before rate decisions. \\
Cost and funding & Direct influence & Align design choices, production rate, and quality controls with executable funding. \\
Process capability and control & Direct influence & Ask whether the process can repeatedly hold the required tolerances at the planned rate. \\
Quality management & Direct influence & Treat prevention and appraisal as cheaper than repeated inspection, rework, and field failure. \\
Manufacturing planning, scheduling, and control & Direct influence & Ensure work instructions, tooling, labor, suppliers, and schedule logic support the required flow. \\
Technology and industrial base & Indirect influence & Assess fragile vendors, unique tooling, special processes, foreign dependencies, and surge capacity. \\
Facilities & Indirect influence & Identify facility constraints, environmental permits, capital equipment, and layout limits. \\
Personnel & Indirect influence & Check whether the supplier has trained labor, certified operators, supervision, and quality culture. \\
\end{longtblr}

\subsection{Manufacturing objectives for design}

The manufacturing objective is to make the design producible, affordable, and
testable without losing the operational capability that justified the program.
The board answer should connect three levers
\autocite{edo-3-7-1-manufacturing-producibility-2025}:

\begin{description}
  \item[\textbf{Design.}] Develop a producible and testable design rather than a design that only works in engineering demonstration.
  \item[\textbf{Process.}] Design and mature manufacturing processes that can reliably build the item.
  \item[\textbf{Process control.}] Establish process controls that satisfy requirements while minimizing avoidable cost.
\end{description}

\subsection{Producibility trade space}

Producibility is a trade, not a slogan.  The program manager should be ready to
explain how operational performance, reliability, availability, supportability,
technology constraints, ownership cost, schedule, and producibility interact.
A design that maximizes one attribute can damage another; for example, a very
tight tolerance may improve a performance margin while increasing supplier
yield loss and delaying delivery
\autocite{edo-3-7-1-manufacturing-producibility-2025}.

\paragraph{Board cue.}
If asked why an EDO cares about manufacturing, answer: production is where
design, funding, industrial base, quality, and schedule converge.  Producibility
must be designed in before production locks in cost and risk.

\end{document}
